% ============================================
% MEM-22 (ZH)
% 中文翻译 — SPEC-935-R356
% ============================================
\section*{\textbf{MEM-22} --- 声音}
\markboth{MEM-22 --- 声音}{}

\noindent \textit{喊叫的病区}

\noindent \textbf{1952 —— 喊叫的病区}

\vspace{0.5em}\noindent\rule{\textwidth}{0.5pt}\vspace{0.5em}

\noindent 躁动病区就在被沉默者病区的隔壁。

\noindent 两个病区。两个世界。一边，是停不下来喊叫的人。另一边，是再也说不出话的人。

\noindent 我待在中间。

\noindent 白天，我听见喊叫。恐惧的喊叫。愤怒的喊叫。没有意义的喊叫。只有声音、没有词语、没有含义的喊叫。

\noindent 夜里，被沉默者的寂静更糟。因为寂静也在喊叫。只是寂静是在你脑子里喊叫。

\noindent 怪物喜欢喊叫。

\noindent \textit{“听，”它说，“每一声喊叫都是一个灵魂在挣脱。每一声喊叫都是某人身上离开的一块。当喊叫停止，那个人已经不在了。”}

\noindent \textit{“不在了，去哪了？”}

\noindent \textit{“进到我里面。”}

\noindent \textit{“你吃喊叫？”}

\noindent \textit{“我什么都吃。吃痛苦。吃恐惧。吃疯狂。都是食物。”}

\noindent 我捂住耳朵。

\noindent 但怪物替我听着。

\vspace{0.5em}\noindent\rule{\textwidth}{0.5pt}\vspace{0.5em}

\noindent 一天夜里，一声不一样的喊叫。

\noindent 不是恐惧。是失去。有人呼喊着某个人的名字。一个女人。她喊着：\textit{“若泽！若泽！”}那名字在走廊里回荡了好几分钟。

\noindent 然后，寂静。

\noindent 第二天早晨，我问护工若泽是谁。

\noindent \textit{“昨天死了。肺炎。”}

\noindent \textit{“那喊叫的是谁？”}

\noindent \textit{“他妻子。1940年两人一起入院的。他死了。她留下了。”}

\noindent 那个女人再也没有喊叫过。

\noindent 她余生在沉默中度过。

\noindent 三个月后，她也死了。

\noindent 人们说是营养不良。

\noindent 我觉得是没有了可以喊叫的人。

\vspace{0.5em}
\noindent\textit{\textbf{↪ 路线:} \rota{MEM-23} \rotasep \rota{MEM-09}}
